\section{Results}
\label{Sec3}

In this section we apply the prescription suggested above and we present our results for the non-perturbative $Z$-factors
both in the RI$'$ and the $\MSb$ schemes. For demonstration purposes we focus on the data with 5 steps of Hypercubic (HYP) smearing that suppress
the power divergence and bring the results closer to renormalized nucleon matrix elements. We focus on the multiplicative renormalization
for the helicity quasi-PDF, and only briefly discuss the case of the unpolarized operator.

\subsection{RI$'$ scheme and conversion to the $\MSb$ scheme}

As a starting point, we have applied the two values of the RI$'$ scale $\bar\mu_0$
used in Fig.~\ref{fig:C} (``parallel'' and ``diagonal''). After converting both cases to the $\MSb$ scheme at $\bar\mu{=}2$ GeV,
we can quantify the systematic uncertainties related to lattice artifacts and truncation of the conversion factor, as explained in the next subsection.

In Fig.~\ref{fig:ZA}, we show the extracted values for the helicity $Z$-factor, $Z_{\Delta h}$, that renormalizes the bare matrix element
$\Delta h(P_3,z)$. In each plot we overlay the results for the RI$'$ (open symbols) and the $\MSb$ (filled symbols) schemes,
for the real and imaginary part of the $Z$-factor. We employ the momentum source technique~\cite{Gockeler:1998ye,Alexandrou:2015sea} that
offers high statistical accuracy with a small number of measurements, typically of ${\cal O}(10)$. As can be seen in the plots, the statistical
uncertainties are almost invisible. The left (right) plot corresponds to the ``parallel'' (``diagonal'') choices for $\bar\mu_0$.
%
We find that the imaginary part of $Z_{\Delta h}^{\overline{\rm MS}}$ is reduced compared to its counterpart in $Z_{\Delta h}^{\rm RI'}$, and is also rather small
for low values of $z$. This is more pronounced in the right plot of Fig.~\ref{fig:ZA} for ``diagonal'' $\bar\mu_0$, for which $\Im[Z_{\Delta h}^{\MSb}]$
(blue circles) is smaller than the corresponding data from the ``parallel'' scale, especially for large values of $z$.
It is worth mentioning that the perturbative $Z$-factor in dimensional regularization and in the $\MSb$ scheme is real to
all orders in perturbation theory, as it is extracted only from the poles. Therefore, it is expected that the imaginary part of the
non-perturbative estimates should be highly suppressed. The behavior of the ``diagonal'' scale is encouraging,
as the imaginary part is very close to zero for $|z|$ up to ${\sim}10\,a$.
\vspace*{0.25cm}

\begin{figure}[h]
\centering
\includegraphics[scale=0.275,angle=-90]{ZA_RI_MS_Landau_pL3_HYP5_mom_parallel.eps}\,\,\,
\includegraphics[scale=0.275,angle=-90]{ZA_RI_MS_Landau_pt_7_pL3_3_3_HYP5_mom3_diagonal.eps}
\vspace*{-0.5cm}
\begin{minipage}{15cm}
\hspace*{3cm}
\caption{\small{The $z$-dependent renormalization function for the matrix element $\Delta h(P_3,z)$
with $a P_3=\frac{6\pi}{L}$. The ``parallel'' and ``diagonal'' choices for $\vec{\bar{\mu}}_0$ are shown in the left and right
plots, respectively. Open (filled) symbols correspond to the RI$'$ ($\MSb$) estimates.}}
\label{fig:ZA}
\end{minipage}
\end{figure}

\vspace*{0.35cm}
As we have done in previous applications of the RI$'$ scheme to extract the renormalization functions of ultra-local operators
(see, e.g., Ref.~\cite{Alexandrou:2015sea}), here we also use a range of values for the RI$'$ renormalization scale,
$\bar\mu_0$. This allows us to study the scale dependence of the vertex functions with the external
momentum and identify the window in  which renormalization factors can be extracted as reliably as possible.
Unlike the case of the local currents, the computation of the lattice artifacts to ${\cal O}(g^2\,a^\infty)$
is extremely laborious and not available yet. Thus, one has to be careful with the choices for $\bar\mu_0$,
as we want to avoid non-perturbative contaminations, and satisfy the criterion that $\hat{P}$ (Eq.~(\ref{Phat})) is small,
preferably below 0.3, to avoid enhanced cut-off effects.
We compute the renormalization functions using the values reported in Table~\ref{tab1}, from which we choose an optimal
range for the fit using the diagonal choices. The parallel $\bar\mu_0$ have only been used to explore the
systematic uncertainties in Subsection~\ref{sub3.2}.

{\small{
\begin{table}[!h]
\vspace*{-0.25cm}
\begin{center}
\begin{tabular}{ccccc}
\hline
\hline\\[-2ex]
Label & $(n_t,n_x,n_y,n_z)$   &  $(a \bar\mu_0)^2$ &  $\bar\mu$ (GeV)  &  $\hat{P}$  \\[0.5ex]
\hline
\hline\\[-2ex]
& & diagonal & & \\
\hline\\[-2ex]
 m1 &(4,3,3,3)		&1.236       &2.671	 &0.261  		\\
 m2 &(5,3,3,3)		&1.332       &2.773	 &0.251 		\\
 m3 &(6,3,3,3)		&1.448       &2.891	 &0.251		\\
 m4 &(7,3,3,3)		&1.583       &3.023	 &0.261 		\\
 m5 &(8,3,3,3)		&1.737       &3.167	 &0.280		\\
 m6 &(9,3,3,3)		&1.911        &3.321	 &0.306 		\\
 m7 &(10,3,3,3)	&2.104       &3.484 	 &0.339		\\
 m8 &(11,3,3,3)	&2.316       &3.656 	 &0.370 		\\[0.5ex]
 \hline\\[-2ex]
 & & parallel & & \\
 \hline\\[-2ex]
 m9 &(4,0,0,3)	&0.542       &1.769 	 &0.539 		\\
 m10 &(9,0,0,3)	&1.216       &2.649 	 &0.592 		\\
 m11 &(11,0,0,3)	&1.622       &3.097 	 &0.664 		\\[0.5ex]
\hline
\hline
\end{tabular}
\vspace*{-0.4cm}
\begin{center}
\begin{minipage}{14cm}
\hspace*{3cm}
\caption{\small{Values for the RI$'$ scale defined as $a \bar\mu_0{=}\frac{2\pi}{L} (\frac{n_t}{2} {+} \frac{\pi}{4}, n_x,n_y,n_z)$,
where $L$ is the spatial extent of the lattice. The values are given in lattice and physical units.
The last column corresponds to the ratio $\hat{P}$ defined in Eq.~(\ref{Phat}).}}
\label{tab1}
\end{minipage}
\end{center}
\end{center}
\vspace*{-0.5cm}
\end{table}
}}

Since we present the renormalized matrix elements for the helicity PDF, we focus on $Z_{\Delta h}$ for this analysis. In Fig.~\ref{fig:Za_vs_scale},
we show $Z^{\MSb}_{\Delta h}$ for selected values of $z$ as a function of the RI$'$ scale, $(a\,\bar\mu_0)^2$, using the
diagonal values labeled by m1$-$m8. The real (imaginary) part is shown in the left (right) panel. We find a residual dependence on $(a\,\bar\mu_0)^2$,
mostly affecting the imaginary part. This is due to the fact that the vertex function depends not only on the magnitude and direction of the renormalization
scale, especially the $z$-direction, as this is parallel to the Wilson line. Upon evolving to the same scale, the $Z$-factors should not depend on the initial scale
if the evolution is known to higher loops in perturbation theory and discretization effects are sufficiently small. The nonzero slope of the plots shown in
Fig.~\ref{fig:Za_vs_scale} is an indication of non-negligible lattice artifacts, as well as, a consequence of the truncation of the conversion factor to one-loop level.
However, further investigation is required in order to attempt disentangling the two effects. We address this in Subsection~\ref{sub3.2}.
\begin{figure}[!h]
\centering
\includegraphics[scale=0.285,angle=-90]{ZA_B55_MS_HYP5_Re.eps}\,\,\,
\includegraphics[scale=0.285,angle=-90]{ZA_B55_MS_HYP5_Im.eps}
\vspace*{-0.3cm}
\begin{minipage}{15cm}
\hspace*{3cm}
\caption{\small{Left: Real part of $Z^\MSb_{\Delta h}$ for scales labeled by m1$-$m8 and $z/a{=}5,10,15$.
Right: Same as left panel for the imaginary part of $Z^\MSb_{\Delta h}$.}}
\label{fig:Za_vs_scale}
\end{minipage}
\end{figure}

\vspace*{0.15cm}
\noindent
To remove the residual dependence on $(a\,\bar\mu_0)^2$ we fit $Z^{\MSb}_{\Delta h} $ with the function
\be
\label{eq:extrapol}
Z^{\MSb}_{\Delta h}  = Z_{0, \,{\Delta h}}^{\MSb}  + Z_{1, \,{\Delta h}}^{\MSb} \, .\,(a\,\bar\mu_0)^2
\ee
and we take $Z_{0, \,{\Delta h}}^{\MSb}$ as our final result. This process is done for each value of
the length of the Wilson line $z$, and repeated for several fit ranges for the diagonal choices m1$-$m8 given in Table~\ref{tab1}.
For the extrapolated value $Z_{0, \,{\Delta h}}^{\MSb}$ at 2 GeV, we choose the one obtained in the range
$(a\,\bar\mu_0)^2\,\epsilon\, [1.4,\,2.0]$. By excluding $(a\,\bar\mu_0)^2$ close to 1, we avoid contamination from non-perturbative
effects. The choice for the above fit range also excludes the two higher scales that have $\hat{P}{>}0.3$, which may carry sizable lattice artifacts.
As an additional check of the quality of the fit, we find that the $\chi^2/d.o.f$ is small for the chosen range.
In Fig.~\ref{fig:ZA_HYP5_extrapolated}, we plot the extrapolated value $Z_{0, \,{\Delta h}}^{\MSb}$ (for 5 steps of HYP smearing), which
is applied on the bare nucleon matrix element of the helicity quasi-PDF in Subsection~\ref{sub3.2}.

\vspace*{0.5cm}
 \begin{figure}[!h]
\centering
\includegraphics[scale=0.275,angle=-90]{ZA_RI_MS_Landau_extrapolated_HYP5_mom3_diagonal.eps}
\vspace*{-0.3cm}
\begin{minipage}{15cm}
\hspace*{3cm}
\caption{\small{Extrapolated $Z$-factor for the helicity operator, $Z_{0, \,{\Delta h}}^{\MSb}$, for 5 steps of HYP smearing.
The employed fit range for the RI$'$ scale is $(a\,\bar\mu_0)^2\,\epsilon\, [1.4 {-} 2.0]$.}}
\label{fig:ZA_HYP5_extrapolated}
\end{minipage}
\end{figure}

\vspace*{0.5cm}
The prescription for obtaining $Z_{VV}$ and $Z_{VS}$ has also been applied on the same ensemble.
We find that both the HYP smearing and the choice of ``diagonal'' scales suppresses the mixing
coefficients $Z_{VS}$ and $Z_{SV}$. In  the left panel of Fig.~\ref{fig:Zvv_Zvs}, we plot
$Z^{\rm RI'}_{VV}$ and $Z^{\rm RI'}_{VS}$ for the scale  ``(7,3,3,3)'' and 5 steps of HYP smearing.
The real and imaginary parts of $Z^{\rm RI'}_{VS}$ are of the same magnitude, but at
least an order of magnitude smaller than the multiplicative factor $Z^{\rm RI'}_{VV}$. This can be
seen from the right panel of Fig.~\ref{fig:Zvv_Zvs}, where we plot the ratios
$\Re[Z^{\rm RI'}_{VS}]/\Re[Z^{\rm RI'}_{VV}]$ (red squares) and
$\Im[Z^{\rm RI'}_{VS}]/\Im[Z^{\rm RI'}_{VV}]$ (blue circles). From the left plot one observes that
the imaginary part of both $Z^{\rm RI'}_{VV}$ and $Z^{\rm RI'}_{VS}$ is compatible with zero for $|z/a|<4$.
Therefore, the large values of the blue points in the region $|z/a|<4$ in the right plot are no indication of significant
mixing.
\begin{figure}[!h]
\centering
\includegraphics[scale=0.285,angle=-90]{Zvv_Zvs_RI_Landau_pt_7_pL3_3_3_HYP5_mom3_diagonal.eps}\,\,\,
\includegraphics[scale=0.285,angle=-90]{Zvs_over_Zvv_RI_Landau_pt_7_pL3_3_3_HYP5_mom3_diagonal.eps}
\vspace*{-0.3cm}
\begin{minipage}{15cm}
\hspace*{3cm}
\caption{\small{Left: multiplicative ($Z^{\rm RI'}_{VV}$) and mixing  ($Z^{\rm RI'}_{VS}$) coefficients entering
the renormalization of the unpolarized quasi-PDF, in the RI$'$ scheme for the ``(7,3,3,3)'' scale. Right:
ratio of the real (red squares) and imaginary (blue circles) parts of $Z^{\rm RI'}_{VS}$ over $Z^{\rm RI'}_{VV}$.}}
\label{fig:Zvv_Zvs}
\end{minipage}
\end{figure}

\subsection{Assessment of systematic uncertainties}
\label{sub3.2}
\medskip

In the framework of this study we have performed several investigations on the systematic uncertainties
related to the truncation of the one-loop conversion factor, as well as, discretization effects. In this subsection
we present the main conclusions of this study, and we give estimates of these effects.

\vspace*{0.3cm}
Prior converting to the $\MSb$ scheme, we compute the $Z$-factors for different values of the RI$'$ scale $\bar\mu_0$.
For example, $Z_{\Delta h}^{\rm RI'}$ shown in the two plots of Fig.~\ref{fig:ZA} on the scales m4 and m9 is different.
Upon evolving to the same scale, the extracted values should agree if the evolution has converged; typically this happens
at two or three loops in perturbation theory. Thus, one should be able to compare the data extracted for
$Z_{\Delta h}^\MSb$ at 2 GeV for the two scales presented in Fig.~\ref{fig:ZA}. The difference between the two, $\Delta Z({\rm m4, m9})$,
is an indication of the presence of lattice artifacts (mainly in the ``parallel'' case), coupled with the truncation of the conversion factor.
This is demonstrated in Fig.~\ref{fig:ZAdiff} for both the real and imaginary parts, and it is interesting to see that the difference
increases as $z$ becomes larger. Based on this observation, we expect that such an increase of the lattice artifacts is also present
for the ``diagonal'' case, but less severe. Another evidence of the presence of non-negligible systematics is the fact that the imaginary
part of $Z_{\Delta h}^\MSb$ is nonzero. In Fig.~\ref{fig:ZA}, we find a small imaginary part for the ``diagonal'' scale, which has an
increasing trend at large values of $z$.

\vspace*{0.3cm}
\begin{figure}[h!]
\centering
\includegraphics[scale=0.3,angle=-90]{Diff_ZMS_parallel_diagonal.eps}
\vspace*{-0.3cm}
\begin{minipage}{15cm}
\hspace*{3cm}
\caption{\small{Difference of $Z_{\Delta h}^\MSb$ between the ``parallel'' (m4) and ``diagonal'' (m9) cases for $\bar\mu_0$.
The real and imaginary parts are shown with red squares and blue circles, respectively.}}
\label{fig:ZAdiff}
\end{minipage}
\end{figure}

\vspace*{0.5cm}
Based on the above, we conclude that the systematic uncertainties related to lattice artifacts and the conversion factor must be addressed,
in order to extract reliable estimates on the renormalization functions. It is our intention to reduce both effects in the near future, which
will eliminate systematic uncertainties propagated to the estimates of the quasi-PDFs.
Understanding the uncertainties dominating the large-$z$ region is crucial, as the matrix element for these values also enters in the Fourier
transform that yields the quasi-PDF. Preliminary explorations indicate that a likely magnitude of the two-loop contribution might suppress
the imaginary part of $Z_{\Delta h}^{\overline{\rm MS}}$. Even though we are currently in no position to accurately quantify these systematics,
we will estimate upper bounds. In particular, we try to estimate the effect of each one individually.


%%%%%%%%%%%%%%%%%%%%%%%%%%%%%%%%%%%%%%%%%%%%%
\vskip 0.75cm
\centerline{\bf\underline{\textit{Lattice artifacts}}}
\vskip 0.25cm
%%%%%%%%%%%%%%%%%%%%%%%%%%%%%%%%%%%%%%%%%%%%%

The $Z$-factor for the helicity quasi-PDF at $z{=}0$ reduces to the renormalization function of the local axial current $Z_A$.
$Z_A$ is scheme- and scale-independent, and thus, the case $z{=}0$ allows us to study the lattice artifacts.  $Z_A$ has been
already evaluated for this ensemble in Ref.~\cite{Alexandrou:2015sea}. In the latter calculation several diagonal scales have
been employed and, together with a technique for the removal of the lattice artifacts, the extracted value was found to be
$Z_A{=}0.7556(5)$. In the present calculation at $z{=}0$ we find a value of $Z_A{=}0.8620(15)$ using the scale m9 and
$Z_A{=}0.7727(2)$ for the scale m4. This is yet another indication that lattice artifacts are large for the  ``parallel''
renormalization scale and less severe, but not negligible, for the ``diagonal'' case.

An additional test that may be performed for any value of $z$ is the comparison of $Z_{\Delta h}^{\rm \MSb}$ between scales
with different components, but same $(a\,\bar\mu_0)^2$. For this, we utilize the values labeled as m10 and m11, which can be compared
to m1 and m4, respectively. Since these pairs are approximately at the same value of $(a\,\bar\mu_0)^2$, we can compare them directly
in the RI$'$ scheme without any evolution\footnote{We confirmed numerically that the conversion factors from m1 to m10, and from m4
to m11 deviates from unity by less than 1$\text{\textperthousand}$, so we ignore it.}, and the difference can be interpreted as lattice artifacts.
We observe similar
behavior for the two pairs and as a demonstration we plot in Fig.~\ref{fig:ZA_diff_m4_m11} the following ratio for each value of $z$
\be
\label{eq:D}
D_{\rm Re\,(Im)}(\bar\mu_0,\bar\mu_0')\equiv \frac{Z_{\rm Re\,(Im)}^{\rm RI'}(\bar\mu_0) - Z_{\rm Re\,(Im)}^{\rm RI'}(\bar\mu_0')}{Z_{\rm Re\,(Im)}^{\rm RI'}(\bar\mu_0)}\,.
\ee
\begin{figure}[h]
\centering
\includegraphics[scale=0.285,angle=-90]{Diff_ZA_B55_ratio_7333vs11300_RI.eps}
\includegraphics[scale=0.285,angle=-90]{Diff_ZA_B55_ratio_5333vs6333_RI.eps}
\vspace*{-0.3cm}
\begin{minipage}{15cm}
\hspace*{3cm}
\caption{\small{Left: The ratio of Eq.~(\ref{eq:D}) using $(a\,\bar\mu_0)^2{=}1.583$ and $(a\,\bar\mu_0')^2{=}1.622$.
Right: Similar to the left plot, for the pairs $( (a\,\bar\mu_0)^2, (a\,\bar\mu_0')^2 ){=}(1.332, 1.448),\, (1.737,1.911)$}}
\label{fig:ZA_diff_m4_m11}
\end{minipage}
\end{figure}

\vspace*{0.25cm}
In the left panel we show the case where $\bar\mu_0$ and $\bar\mu_0'$ correspond to the values m4 and m11, respectively.
The difference from zero can be attributed to lattice artifacts. We find that the $Z$-factor extracted from m11 deviates from the value extracted
from m4 by 5-10\% in the real part and up to 30\% in the imaginary part. This deviation seems to stabilize after $z/a{\sim}8$ to 5\% (10\%)
for the real (imaginary) part.

Note that the left plot of Fig.~\ref{fig:ZA_diff_m4_m11} gives the estimate of the discretization effects for the ``parallel''
case compared to the ``diagonal'' one. It would be interesting to compute $D_{\rm Re\,(Im)}$ between two neighboring ``diagonal''
momenta in order to understand the change in the artifacts. We choose 2 pairs (m2,m3) and (m5,m6) and we find that both
$D_{\rm Re}(\rm{m2,m3})$ and $D_{\rm Re}(\rm{m5,m6})$ are less than 1\%, while $D_{\rm Im}(\rm{m2,m3})$ and $D_{\rm Im}(\rm{m5,m6})$
are of the order of 6\%. This is a confirmation that excluding ``parallel'' momenta in the fit of Eq.~(\ref{eq:extrapol}) reduces
significantly the discretization effects.

%%%%%%%%%%%%%%%%%%%%%%%%%%%%%%%%%%%%%%%%%%%%%
%\vskip 0.75cm
\newpage

\centerline{\bf\underline{\textit{Truncation effects}}}
\vskip 0.25cm
%%%%%%%%%%%%%%%%%%%%%%%%%%%%%%%%%%%%%%%%%%%%%

In order to assess quantitatively the influence of the conversion factor that is known to one-loop perturbation theory, we form the ratios
\be
\label{eq:R}
R_{\rm Re\,(Im)}^{{\rm RI'}\,(\MSb)}(z, \bar\mu_0,\bar\mu_0';\bar\mu)\equiv
\frac{Z_{\rm Re\,(Im)}^{{\rm RI'}\,(\MSb)}(z,\bar\mu_0;\bar\mu)}{Z_{\rm Re\,(Im)}^{{\rm RI'}\,(\MSb)}(z,\bar\mu_0';\bar\mu)}
\ee
both for the real and imaginary parts of the helicity $Z$-factor.  In Fig.~\ref{fig:RIvsMS}, we plot Eq.~(\ref{eq:R}) for the RI' and
$\MSb$ case. These have been extracted at different values of $(a\bar\mu_0)^2$ (m1$-$m8) and evolved perturbatively to the same
scale of $2$ GeV, using the results of Ref.~\cite{Constantinou:2017sej}. The ratio is always taken with respect to the
smallest ``diagonal'' scale, m1. $R_{\rm Re\,(Im)}^{\rm RI'}$ depend on the truncation effects in the one-loop evolution to 2 GeV, while
$R_{\rm Re\,(Im)}^\MSb$ is affected by scheme conversion truncation effects. Without contamination from lattice artifacts and
truncation effects, Eq.~(\ref{eq:R}) should equal 1 in both schemes and for all values of $(a\bar\mu_0)^2$ and $z/a$. This realization
allows for investigation of the truncation effects.

\bigskip
\begin{figure}[h]
\centering
\includegraphics[scale=0.285,angle=-90]{ZA_B55_RI2GEV_MS_HYP5_Re.eps}\,\,\,
\includegraphics[scale=0.285,angle=-90]{ZA_B55_RI2GEV_MS_HYP5_Im.eps}
\vspace*{-0.3cm}
\begin{minipage}{15cm}
\hspace*{3cm}
\caption{\small{The ratio of Eq.~(\ref{eq:R}) for $z/a{=}5,10,15$ as a function of the initial RI$'$ scale.
All results have been evolved to 2 GeV and are normalized with the value of the $Z$-factor at
$(a\bar\mu_0')^2{=}1.236$ (m1). Left: real part. Right: imaginary part. }}
\label{fig:RIvsMS}
\end{minipage}
\end{figure}

One observes that the imaginary part is more sensitive to the change of the initial RI$'$ scale, which is given in the x-axis.
To demonstrate this, we keep the range of y-axis the same for all plots of Fig.~\ref{fig:RIvsMS}.
The difference between the ratios of the real part extracted from different RI$'$ scales $(\bar\mu_0)^2$ is consistently
small and reaches at most approx.\ 5\% in the RI$'$ scheme and 3\% in the $\MSb$ scheme, regardless of the Wilson line
length, $z$.  As we mentioned above, for $z{=}0$ we can compare to the highly reliable value found in Ref.~\cite{Alexandrou:2015sea};
our current value from the scale m1 is around 2\% higher. Hence, we can estimate the typical size of
discretization effects to be of order 2-5\% in the real part of the $Z$-factors.

The differences in the $\MSb$ ratios are combinations of lattice artifacts and the conversion truncation.
We observe that the conversion to the $\MSb$ scheme decreases the differences in the $\MSb$ scheme to at most 3\%.
Hence, truncation alone in the real part seems not to have an effect larger than 2\%. Since we know that conversion
truncation effects are very small for small values of $z$ (in particular, they are zero for $z{=}0$, where $Z_A$
is scheme-independent), we can take 0-2\% as our estimate. In the end, given the fact that truncation
effects seem to be opposite to the influence of lattice artifacts, we estimate that the total effects present in the
real part of $Z_{\Delta h}^\MSb$ should not exceed 5\%.

The situation is somewhat different in the ratio of the imaginary parts shown in the right panel of Fig.~\ref{fig:RIvsMS}.
We conclude that the uncertainty from lattice artifacts can be of the order of 10\% in the RI$'$ scheme for intermediate
and large $z$ ($|z|/a\geq5$\,\footnote{For smaller values of $z$, the imaginary part is small in absolute terms and
Eq.~(\ref{eq:R})  becomes meaningless, i.e.\ even small absolute changes of the imaginary part of $Z_{\Delta h}^\MSb$
can imply large changes in $R$.}). From the right panel of Fig.~\ref{fig:RIvsMS}, the total uncertainty in
$R_{\rm Im}^\MSb$ can get enhanced to around 40\% for large and up to 80\% for intermediate Wilson line lengths.
In addition, the total uncertainty in the imaginary part of $Z_{\Delta h}^\MSb$ is basically of the order of its magnitude,
as it is expected to be real. Note that ignoring the imaginary part does not provide any solution, as
 they need to come out zero from the computation in order to claim that the computation is fully reliable.
Thus, conversion truncation effects become the most major source of uncertainty of the imaginary parts of the $Z$-factors.

\bigskip
Based on the study presented in this subsection, we present in Table~\ref{tab2} a summary of our quantitative estimates on
the systematic uncertainties present in $Z_{\Delta h}^\MSb$. The estimate of the total effect is not a simple sum of
isolated effects, but takes into account their different signs discussed above. For the imaginary part, the estimates are for
$|z/a|\geq5$, since $\Im[Z_{\Delta h}^\MSb]{\sim}0$ for smaller values of $z$ and the relative effects become meaningless.

We would like to stress that due to the complex multiplication of the $Z$-factors and the bare matrix elements,
the large uncertainty in the imaginary part of the $Z$-factor implies that also real part of the renormalized matrix
elements is affected. Furthermore, the uncertainties that we consider in the $Z$-factors translate differently to the
uncertainties of the renormalized matrix elements for different Wilson line lengths and thus, we postpone the discussion of this
influence to the next subsection.

{\small{
\begin{table}[!h]
\vspace*{0.25cm}
\begin{center}
\begin{tabular}{ccc}
\hline
\hline\\[-2ex]
Effect & $\Re[Z_{\Delta h}^\MSb]$   &  $\Im[Z_{\Delta h}^\MSb]$ \\[0.5ex]
\hline
\hline\\[-2ex]
Lattice artifacts & 2-5\% & $\lesssim10$\% \\[1ex]
Evolution truncation& 1-2\% & 1-2\% \\[1ex]
Conversion truncation& $\lesssim2$\% & $\lesssim100$\% \\[2ex]
\hline\\[-1.5ex]
Total & 3-5\% &  $\lesssim100$\% \\[1ex]
\hline
\hline
\end{tabular}
\vspace*{-0.4cm}
\begin{center}
\begin{minipage}{14cm}
\hspace*{3cm}
\caption{\small{Quantitative estimates of systematic uncertainties in the real and imaginary parts of $Z_{\Delta h}^\MSb$.}}
\label{tab2}
\end{minipage}
\end{center}
\end{center}
\end{table}
}}

\subsection{Renormalized results}
\label{sub3.2}

Once the $Z$-factors are obtained for the unpolarized, helicity and transversity quasi-PDFs, one may
proceed with the application of the renormalization in the nucleon matrix elements. Here we mostly focus
on the helicity case, as the renormalization is multiplicative. The case of the unpolarized quasi-PDF is
briefly mentioned in the end of the subsection.

In Fig.~\ref{fig:deltah}, we show the renormalized helicity nucleon matrix elements, and compare them with
the bare ones. This is a straightforward procedure as there is no mixing for the axial operator and, therefore,
the renormalization is only multiplicative:
\be
\Delta h^{\MSb} (z) = Z_{\Delta h}^\MSb(z)\cdot \Delta h^{bare}(z)\,\,.
\ee
The above formula involves complex quantities, and thus, $\Delta h^{\MSb} (z)$ is a mixture of the real and imaginary part
of the bare matrix element. However, each value of $z$ is renormalized independently.
We use $Z_{\Delta h}^\MSb$ from the extrapolation of Eq.~(\ref{eq:extrapol}) in the range
$(a\,\bar\mu_0)^2\,\epsilon\, [1.4,\,2.0]$, as shown in Fig.~\ref{fig:ZA_HYP5_extrapolated}.
One can see in Fig.~\ref{fig:deltah} that for small values of $z$ there is a slight suppression of
the renormalized real part with respect to the bare one ($\Re [Z_{\Delta h}^\MSb<1]$).
$\Re[\Delta h^{\MSb}]$ is compatible with zero for $|z/a|{>}8$, but with increased statistical
uncertainties. The effect of the renormalization on the imaginary part of the matrix element is
profound in the large $z$ region, where we observe an amplification of its value and a shift of the
maximum to larger z, as compared to the bare one.
\bigskip
\begin{figure}[h]
\centering
\includegraphics[scale=0.285,angle=-90]{deltah_renorm_pL3_HYP5_Re.eps}\,\,\,
\includegraphics[scale=0.285,angle=-90]{deltah_renorm_pL3_HYP5_Im.eps}
\vspace*{-0.3cm}
\begin{minipage}{15cm}
\hspace*{3cm}
\caption{\small{Renormalized matrix elements for the helicity quasi-PDF in the
$\MSb$ scheme at $\bar\mu{=}$2 GeV using $Z_{0, \Delta h}^{\MSb}$
extracted from the fit range $(a\,\bar\mu_0)^2\,\epsilon\, [1.4,\,2.0]$.}}
\label{fig:deltah}
\end{minipage}
\end{figure}

\vspace*{0.5cm}
The renormalized matrix elements of the helicity quasi-PDF are presented here for the first time, and they demonstrate the
enormous progress in the field of the quasi-PDFs. However, before attempting to compare with the physical PDFs, we must
understand the uncertainties that are inherited to $\Delta h^{\MSb} (z) $ from its renormalization function.
As we argued in the previous subsection, a robust computation needs the subtraction of $\mathcal{O}(g^2\,a^\infty)$
lattice artifacts and a significant reduction of truncation uncertainties in the conversion between the RI$'$ and $\MSb$ schemes.
We attempt setting bounds on these systematics, starting with the real part of renormalized matrix elements, which reads:
\be
\label{eq:Re}
\Re[\Delta h^\MSb]=\Re[Z_{\Delta h}^{\MSb}]\,\, \Re[\Delta h^{bare}]-\Im[Z_{\Delta h}^{\MSb}]\,\, \Im[\Delta h^{bare}]\,.
\ee
For small values of $z$, $\Im[Z_{\Delta h}^{\MSb}]$ is approximately zero and thus $\Re[\Delta h^\MSb]$ is approximately equal to the first term of Eq.~(\ref{eq:Re}).
On the contrary, for $|z/a|\gtrsim10$, $\Re[\Delta h^\MSb]$ receives significant contributions from the imaginary part, because $\Re[\Delta h^{bare}]<\Im[\Delta h^{bare}]$.


The uncertainty in the small-$z$ region of renormalized matrix elements is dominated by uncertainties in the real part of the $Z$-factor, which, as argued in the previous subsection, should not exceed 5\%.
The local minimum observed at $z{=}0$ is likely to result from lattice artifacts and truncation effects.
In the large-$z$ region, the real part of the renormalized matrix elements receives negative contributions from the imaginary part of the bare matrix element if $Z_{\Delta h}^{\MSb}$ has a non-zero imaginary part and $\Im[\Delta h^{bare}]$ decays to zero more slowly than $\Re[\Delta h^{bare}]$, which is what we observe in the data.
This effect is unphysical and is expected to be strongly suppressed or eliminated by extending our calculation to the two-loop order in the conversion and by subtracting lattice artifacts.

The imaginary part of renormalized matrix elements,
\be
\label{eq:Im}
\Im[\Delta h^\MSb]=\Re[Z_{\Delta h}^{\MSb}]\,\, \Im[\Delta h^{bare}]+\Im[Z_{\Delta h}^{\MSb}]\,\, \Re[\Delta h^{bare}],
\ee
is enhanced in the intermediate-$z$ regime, compared to the bare matrix elements. This results mostly from the fact that
$\Re[Z_{\Delta h}^{\MSb}]$ increases with increasing $z$ at a faster rate, than the decay of $\Im[\Delta h^{bare}]$.
However, the obtained values receive also contributions from the second term of Eq.~(\ref{eq:Im}), where $\Im[Z_{\Delta h}^{\MSb}]$
is subject to a large uncertainty, as discussed in the previous subsection.

We want to stress that it is not possible to give a single number for the relative uncertainty of the real and imaginary parts of $\Delta h^\MSb$.
According to Eqs.~(\ref{eq:Re}) - (\ref{eq:Im}), different regions of $z$ are influenced in a different way by the real and imaginary parts of the $Z$-factors.
For small $z$, where $\Im[Z_{\Delta h}^{\MSb}]$ is small, the propagated uncertainty in the matrix elements is dominated by the uncertainty of
 $\Re[Z_{\Delta h}^{\MSb}]$, which is of the order of 5\% and comparable to the currently attained statistical uncertainty. Thus, $\Delta h^\MSb$
is rather robust in this region. However, when $z$ is increased, the uncertainty from $\Im[Z_{\Delta h}^{\MSb}]$ starts to dominate and reaches
100\% for large $z$. In this way, the improvements expected by the perturbative subtraction of $\mathcal{O}(g^2\,a^\infty)$ lattice artifacts and the extension
to the two-loop perturbative conversion to the $\MSb$ scheme are very important for obtaining meaningful values of renormalized matrix elements and hence, also quasi-PDFs.

\vspace*{0.5cm}
The values of the multiplicative vector renormalization factor and the mixing coefficient can be used to properly
renormalize the unpolarized quasi-PDF through Eq.~(\ref{h_R}). The successful renormalization requires the bare
nucleon matrix elements for the scalar and vector operators, in order to extract the renormalized $h_V^{\rm RI'}$.
Then, it must be multiplied by the conversion factor $C_V$ to bring the results to the $\MSb$ scheme.
Note that once the mixing is treated, the conversion factor is multiplicative and not a $2{\times}2$ matrix.
This is due to the fact that no mixing is present in the continuum dimensional regularization.


\subsection{Matching to light-front PDFs}
Having the renormalized matrix elements, one can perform a Fourier transform and obtain the renormalized quasi-PDF,
which represents the distribution of quark momenta for a finite-momentum nucleon moving in the chosen spatial ($z$) direction:
\begin{equation}
\label{qpdf}
\tilde{q}(x,\mu,P_3) = \int_{-\infty}^\infty \frac{dz}{4\pi} e^{-i z x P_3} \langle N |
\bar{\psi}(0,z)\Gamma W(z) \psi(0,0) |N\rangle^{\MSb,\mu}\,,
\end{equation}
where $x$ is the quark momentum fraction.
The quasi-PDF, expressed in our case in the $\MSb$ scheme at $\mu{=}2$ GeV, can then be connected to the light-front PDF, in the same scheme and at the same scale, using one-loop perturbative matching.
The matching procedure uses the fact that only the ultraviolet physics is different in quasi- ($\tilde{q}$) and light-front ($q$) PDFs \cite{Xiong:2013bka}.
Hence, the one-loop difference between them is expressed as the difference between vertex corrections (denoted $Z^{(1)}$ below) and wave function corrections ($\delta Z^{(1)}$) for the finite momentum and infinite momentum cases.
The generic matching formula for quasi-PDFs is: \cite{Ji:2015jwa,Chen:2016fxx}
\begin{eqnarray}
\label{invq}
q\left(x,\mu\right) &=& \tilde{q}(x,\mu,P_3) - \frac{\alpha_s(\mu)}{2\pi} \,\tilde{q}(x,\mu,P_3)\,\delta Z^{(1)}\!\left( \frac{\mu}{P_3}\right)  \nonumber\\*
& & - \frac{\alpha_s(\mu)}{2\pi} \int_{-\infty}^{\infty} Z^{(1)}\!\left( \xi, \frac{\mu}{P_3} \right) \tilde{q}\! \left( \frac{x}{\xi},\mu,P_3 \right) \frac{d\xi}{|\xi |} + \mathcal{O}(\alpha_s^2),
\end{eqnarray}
where $\alpha_s(\mu)$ is the strong coupling constant at the scale $\mu$.
In the integral, we exclude a small region around $\xi=0$, such that the argument of $\tilde{q}$ is not too large, as we would then pick up contributions from the mirror images of $\tilde{q}$ resulting from the periodicity of the Fourier transform in Eq.~(\ref{qpdf}).
The linearly divergent terms $\propto\Lambda/P_3$ ($\Lambda$: transverse momentum cutoff) of the matching
formulae from Ref.~\cite{Xiong:2013bka} are not present in our renormalized results.

\bigskip
\begin{figure}[h]
\centering
\includegraphics[scale=0.55,angle=-90]{qmatched_bare_vs_fit69v2.eps}\,\,\,
\vspace*{-0.3cm}
\begin{minipage}{15cm}
\hspace*{3cm}
\caption{\small{Comparison of matched helicity PDF obtained from quasi-PDF computed with either fully renormalized matrix elements (blue)
or with bare matrix elements multiplied by the local ($z{=}0$) axial current $Z$-factor, $Z_A$ (magenta). For purely orientational purposes,
we also plot phenomenological PDFs (DSSV08 \cite{deFlorian:2009vb} and JAM15 \cite{Sato:2016tuz}).
However, we emphasize that no quantitative comparison with our results is aimed at, since careful consideration of a number of systematic effects is still needed. These include:
cut-off effects, non-physical pion mass, finite volume effects, possible contamination by excited states, extrapolation to infinite nucleon boost, as well as the improvements in the computation of $\MSb$ renormalization functions, postulated in the previous subsection.}}
\label{fig:matched}
\end{minipage}
\end{figure}

\medskip
In Fig.~\ref{fig:matched}, we show the matched helicity PDF computed with either fully renormalized matrix elements obtained in this work (blue) or with bare matrix elements multiplied by the local ($z{=}0$) axial vector current renormalization function $Z_A$ (magenta), corresponding to our results from Ref.~\cite{Alexandrou:2016jqi}.
We observe that the renormalized matrix elements from this work move towards the phenomenological PDFs, which is promising.
In particular, the antiquark asymmetry is not overestimated any longer and actually this asymmetry becomes compatible with zero under current uncertainties.
In the quark part, there is an enhancement of the matched PDF for all values of $x$.
We emphasize again that the comparison with the phenomenological PDFs should be understood as only qualitative.
For quantitative comparison, a careful investigation of a number of systematic effects is still needed. These include:
cut-off effects, non-physical pion mass, finite volume effects, possible contamination by excited states, extrapolation
to infinite nucleon boost, as well as the improvements in the computation of $\MSb$ renormalization functions,
postulated in the previous subsection: subtraction of lattice artifacts computed in lattice perturbation theory and
reduction of truncation effects in the perturbative conversion to the $\MSb$ scheme.
